% ================ 第八节 公司治理与独立性 ====================================
% 治理制度与股东会/董事会运作（新《公司法》：不设监事会，审计委员会行使
% 监事会职权）→ 独立董事与专门委员会 → 内部控制 → 独立性（五独立）
% → 同业竞争 → 关联方与关联交易 → 违规担保与资金占用。
\gssection{公司治理与独立性}

\subsection{公司治理制度的建立健全及运行情况}

公司已按照《公司法》《证券法》及中国证监会、上交所的有关规定，建立了由股东
会、董事会和高级管理层构成的公司治理架构，制定了《公司章程》《股东会议事规
则》《董事会议事规则》《独立董事工作制度》《对外担保管理制度》《关联交易管
理制度》《对外投资管理制度》等治理制度。根据《公司法》及《公司章程》的规定，
公司不设监事会，由董事会审计委员会行使《公司法》规定的监事会职权。

报告期内，公司股东会、董事会均依照法律法规及公司制度规范召集、召开，会议记
录完整，决策程序合法有效，不存在损害公司及股东合法权益的情形。

\subsection{独立董事及董事会专门委员会}

公司董事会由9名董事组成，其中独立董事3名（含会计专业人士1名）、职工代表董
事1名。董事会下设审计委员会、提名委员会、薪酬与考核委员会、战略委员会四个
专门委员会：审计委员会由独立董事赵××、钱××、孙××组成，赵××（会计专
业人士）任召集人，并行使《公司法》规定的监事会职权；提名委员会、薪酬与考核
委员会均由独立董事占多数并担任召集人。报告期内，独立董事及各专门委员会依照
相关制度履行职责，对公司重大事项发表了独立、客观的意见。

\subsection{内部控制情况}

公司已建立与经营规模相适应的内部控制体系，覆盖销售与收款、采购与付款、项目
交付、研发与数据管理、资金、关联交易、对外担保、信息披露等主要业务环节。董
事会对公司内部控制进行了自我评价，认为公司内部控制健全且被有效执行。示例会
计师对公司内部控制出具了《内部控制鉴证报告》（示例会所鉴字〔2026〕第××号），
认为公司于2025年12月31日在所有重大方面保持了与财务报表相关的有效的内部控制。

\subsection{发行人的独立性}

\subsubsection{资产完整}

公司拥有与生产经营有关的房产、设备、专利、软件著作权、商标及域名等资产的完
整所有权或使用权，资产独立完整，不存在被控股股东、实际控制人及其控制的其他
企业占用的情形。

\subsubsection{人员独立}

公司总经理、副总经理、财务总监、董事会秘书等高级管理人员均专职在公司工作并
领取薪酬，未在控股股东、实际控制人及其控制的其他企业中担任除董事以外的其他
职务；公司财务人员未在上述企业中兼职。公司建立了独立的劳动、人事及薪酬管理
体系。

\subsubsection{财务独立}

公司设立了独立的财务部门，建立了独立的财务核算体系和财务管理制度，独立开设
银行账户，依法独立纳税，独立作出财务决策，不存在与控股股东、实际控制人共用
银行账户的情形。

\subsubsection{机构独立}

公司建立了健全的内部经营管理机构，独立行使经营管理职权，与控股股东、实际控
制人及其控制的其他企业间不存在机构混同的情形。

\subsubsection{业务独立}

公司独立开展研发、采购、交付与销售，拥有独立完整的业务体系及直接面向市场独
立经营的能力，不存在对控股股东、实际控制人及其控制的其他企业的业务依赖。

\subsection{同业竞争}

截至本招股说明书签署日，控股股东示例控股为持股型公司，实际控制人林××控制
的其他企业均不从事与公司相同或相似的业务，公司与控股股东、实际控制人及其控
制的其他企业之间不存在同业竞争。

为避免未来可能发生的同业竞争，控股股东示例控股及实际控制人林××已分别出具
《关于避免同业竞争的承诺函》，承诺其及其控制的其他企业不以任何形式直接或间
接从事与公司主营业务构成竞争的业务。

\subsection{关联方及关联交易}

\subsubsection{主要关联方}

公司主要关联方包括：控股股东示例控股、实际控制人林××及其近亲属、员工持股
平台示例聚力、持股5\%以上的股东示例创投与示例产投、公司董事及高级管理人员
及其关系密切的家庭成员，以及上述主体控制或施加重大影响的其他企业。

\subsubsection{经常性关联交易}

报告期内，公司不存在采购、销售商品或者接受、提供劳务等经常性关联交易。

\subsubsection{偶发性关联交易}

2023年度，公司向示例控股租赁办公场所，租赁费用为128.45万元，定价参考同地段
市场租金水平，作价公允；该租赁已于2023年12月到期终止。除此之外，报告期内公
司无其他偶发性关联交易。

\subsubsection{关联交易的决策程序及独立董事意见}

公司已在《公司章程》及《关联交易管理制度》中明确关联交易的审议权限、回避表
决及定价原则。独立董事对报告期内关联交易发表意见认为：报告期内的关联交易定
价公允，履行了必要的决策程序，不存在损害公司及其他股东利益的情形。

\subsection{违规担保及资金占用情况}

报告期内，公司不存在为控股股东、实际控制人及其控制的其他企业进行违规担保的
情形，不存在资金被控股股东、实际控制人及其关联方违规占用的情形。
